\chapter{The Born Measure and the Universal Prior}\label{app:born}

Chapter 8 compared the Born rule and the universal prior without identifying them. This appendix fixes the boundary of that analogy.

Start with what each weighting does. Quantum theory uses the Born rule to assign probabilities from squared amplitudes $|\psi|^2$. In an Everettian interpretation, those weights are associated with decoherent outcomes, but unitary evolution and decoherence alone do not prove that outcomes are ontically separate worlds. Algorithmic information theory fixes a universal machine and assigns fair-bit weights $2^{-|p|}$ to programs, which are summed into $M$ over common output prefixes. The Born rule weights quantum outcomes; $M$ weights program descriptions. Neither formula by itself creates or counts observer occurrences.

The analogy is therefore limited to a role: both attach quantitative weight after a domain of alternatives has been specified. Their domains, normalization, empirical status, and mathematical justifications differ.

The book does not derive the Born rule from $M$. No calculation here starts from algorithmic information theory and yields $|\psi|^2$. Nor is flat branch-counting a ready alternative: decoherence does not define a unique number of branches. The conceptual and mathematical work of justifying Born statistics within a quantum interpretation remains open to the relevant quantum theory.

The postulate also does not make that work disappear. Denying a view-independently privileged decomposition prevents a representation from certifying itself, but experiments still exhibit Born statistics and a quantum theory still owes an account of them. ``Why this weighting?'' can be an ordinary scientific and mathematical question without presupposing a view from nowhere.

This is the seam to \textit{Multitudes}, the companion volume that develops an Everettian interpretation. This book borrows the landmark and hands back a caution: program weights neither derive quantum weights nor settle their ontology. \textit{Multitudes} need not be read first; the distinction above is the whole of what this appendix requires.
